% GENERATED FILE. Edit canonical lessons, then run npm run generate:books.
% canonical-source-hash: fnv1a64:26640e2b92b8728d
% canonical-lessons: MR-C492-ucki-lagne, MR-C492-dhekar-dene, MR-C492-thunkne, MR-C492-dhap-lagne, MR-C492-thaska-lagne

\chapter{Hiccups and Gasps}
\label{ch:mr-hiccups-and-gasps}
\hlchaptermodality{492}

\begin{chapteropening}
\textbf{By the end of this chapter:} \emph{I can talk about hiccups, burps, spitting and being out of breath.}

Complete the last lesson of chapter 492: I can talk about hiccups, burps, spitting and being out of breath.
\end{chapteropening}

\section[\texorpdfstring{\mr{उचकी} \mr{लागणे} (uckī lāgṇe)}{उचकी लागणे (uckī lāgṇe)}]{\mr{उचकी} \mr{लागणे} (uckī lāgṇe) — to have hiccups}

\label{lesson:MR-C492-ucki-lagne}



\begin{quote}
Before the new one: say the Marathi for to take a nap, then the Marathi for to lie down.
\end{quote}

\subsection*{You'll want to know: \mr{उचकी} \mr{लागणे}}
\textbf{\mr{उचकी} \mr{लागणे}} — \emph{uckī lāgṇe} — \textquotedblleft{}to have hiccups\textquotedblright{}.

\begin{quote}
\mr{घाईत} \mr{जेवल्याने} \mr{त्याला} \mr{उचकी} \mr{लागली}.
\end{quote}

\begin{quote}
\emph{ghāīt jevlyāne tyālā uckī lāglī.}
\end{quote}

\begin{quote}
Eating in a hurry gave him hiccups.
\end{quote}

Past (to him …): \mr{त्याला} \mr{उचकी} \mr{लागली}, \emph{tyālā uckī lāglī}.

Dative: \mr{मला} \mr{उचकी} \mr{लागली}.

\begin{cousinweb}
\textbf{\mr{उचकी} \mr{लागणे}} (\emph{uckī lāgṇe}) is a Marathi phrase built on \mr{उचकी}, \textquotedblleft{}a hiccup\textquotedblright{}, whose origin is not settled, with \mr{लागणे}.
\end{cousinweb}

\begin{grammarlens}[title={What you've built}]
One more word for describing what you lived and felt.
\end{grammarlens}

\subsection*{Your turn}
\begin{itemize}
\raggedright
  \item \emph{Say it:} \emph{uckī lāgṇe}
  \item \emph{Say it:} \emph{uckī lāgṇe}, once more
  \item \emph{Say it:} \emph{ghāīt jevlyāne tyālā uckī lāglī.} — Eating in a hurry gave him hiccups.
  \item \emph{Recall:} say the Marathi for to take a nap, then the Marathi for to lie down, then say \emph{uckī lāgṇe} again
\end{itemize}

\begin{tcolorbox}[breakable,colback=teal!4,colframe=teal!35!black,title={Before you move on}]
What does \mr{उचकी} \mr{लागणे} mean? (\textquotedblleft{}To have hiccups\textquotedblright{}.) Say it once more.
\end{tcolorbox}
\section[\texorpdfstring{\mr{ढेकर} \mr{देणे} (ḍhekar deṇe)}{ढेकर देणे (ḍhekar deṇe)}]{\mr{ढेकर} \mr{देणे} (ḍhekar deṇe) — to burp}

\label{lesson:MR-C492-dhekar-dene}



\begin{quote}
Before the new one: say the Marathi for to lie down, then the Marathi for to have hiccups.
\end{quote}

\subsection*{You'll want to know: \mr{ढेकर} \mr{देणे}}
\textbf{\mr{ढेकर} \mr{देणे}} — \emph{ḍhekar deṇe} — \textquotedblleft{}to burp\textquotedblright{}.

\begin{quote}
\mr{पोटभर} \mr{जेवल्यावर} \mr{आजोबांनी} \mr{ढेकर} \mr{दिली}.
\end{quote}

\begin{quote}
\emph{poṭbhar jevlyāvar ājobānnī ḍhekar dilī.}
\end{quote}

\begin{quote}
After a full meal Grandfather burped.
\end{quote}

Past, with -\mr{ने} (he …): \mr{त्याने} \mr{ढेकर} \mr{दिली}, \emph{tyāne ḍhekar dilī}.

\mr{ढेकर}, a burp, is feminine in everyday use (\mr{ढेकर} \mr{आली}, \mr{ढेकर} \mr{दिली}), though CDIAL marks it masculine. \mr{ढेकर} \mr{येणे}, with the dative.

\begin{cousinweb}
\textbf{\mr{ढेकर} \mr{देणे}} (\emph{ḍhekar deṇe}) is a Marathi phrase built on \mr{ढेकर}, \textquotedblleft{}a burp\textquotedblright{}, which is inherited: CDIAL lists Marathi \emph{ḍhẽkar}, \textquotedblleft{}belch\textquotedblright{}, under \emph{*ḍakkāra} (CDIAL 5521); \mr{देणे} is the light verb.
\end{cousinweb}

\begin{grammarlens}[title={What you've built}]
One more word for describing what you lived and felt.
\end{grammarlens}

\subsection*{Your turn}
\begin{itemize}
\raggedright
  \item \emph{Say it:} \emph{ḍhekar deṇe}
  \item \emph{Say it:} \emph{ḍhekar deṇe}, once more
  \item \emph{Say it:} \emph{poṭbhar jevlyāvar ājobānnī ḍhekar dilī.} — After a full meal Grandfather burped.
  \item \emph{Recall:} say the Marathi for to lie down, then the Marathi for to have hiccups, then say \emph{ḍhekar deṇe} again
\end{itemize}

\begin{tcolorbox}[breakable,colback=teal!4,colframe=teal!35!black,title={Before you move on}]
What does \mr{ढेकर} \mr{देणे} mean? (\textquotedblleft{}To burp\textquotedblright{}.) Say it once more.
\end{tcolorbox}
\section[\texorpdfstring{\mr{थुंकणे} (thuṅkṇe)}{थुंकणे (thuṅkṇe)}]{\mr{थुंकणे} (thuṅkṇe) — to spit}

\label{lesson:MR-C492-thunkne}



\begin{quote}
Before the new one: say the Marathi for to have hiccups, then the Marathi for to burp.
\end{quote}

\subsection*{You'll want to know: \mr{थुंकणे}}
\textbf{\mr{थुंकणे}} — \emph{thuṅkṇe} — \textquotedblleft{}to spit\textquotedblright{}.

\begin{quote}
\mr{रस्त्यावर} \mr{थुंकणं} \mr{चुकीचं} \mr{आहे}.
\end{quote}

\begin{quote}
\emph{rastyāvar thuṅkṇaṁ cukīcaṁ āhe.}
\end{quote}

\begin{quote}
Spitting in the street is wrong.
\end{quote}

Past (he …): \mr{तो} \mr{थुंकला}, \emph{to thuṅklā}.

\mr{थुंकी}, spittle.

\begin{cousinweb}
\textbf{\mr{थुंकणे}} (\emph{thuṅkṇe}) is inherited: CDIAL lists Marathi \emph{thũkṇẽ}, \textquotedblleft{}spits\textquotedblright{}, under the expressive \emph{*thukk} (CDIAL 6097).
\end{cousinweb}

\begin{grammarlens}[title={What you've built}]
One more word for describing what you lived and felt.
\end{grammarlens}

\subsection*{Your turn}
\begin{itemize}
\raggedright
  \item \emph{Say it:} \emph{thuṅkṇe}
  \item \emph{Say it:} \emph{thuṅkṇe}, once more
  \item \emph{Say it:} \emph{rastyāvar thuṅkṇaṁ cukīcaṁ āhe.} — Spitting in the street is wrong.
  \item \emph{Recall:} say the Marathi for to have hiccups, then the Marathi for to burp, then say \emph{thuṅkṇe} again
\end{itemize}

\begin{tcolorbox}[breakable,colback=teal!4,colframe=teal!35!black,title={Before you move on}]
What does \mr{थुंकणे} mean? (\textquotedblleft{}To spit\textquotedblright{}.) Say it once more.
\end{tcolorbox}
\section[\texorpdfstring{\mr{धाप} \mr{लागणे} (dhāp lāgṇe)}{धाप लागणे (dhāp lāgṇe)}]{\mr{धाप} \mr{लागणे} (dhāp lāgṇe) — to get breathless}

\label{lesson:MR-C492-dhap-lagne}



\begin{quote}
Before the new one: say the Marathi for to burp, then the Marathi for to spit.
\end{quote}

\subsection*{You'll want to know: \mr{धाप} \mr{लागणे}}
\textbf{\mr{धाप} \mr{लागणे}} — \emph{dhāp lāgṇe} — \textquotedblleft{}to get breathless\textquotedblright{}.

\begin{quote}
\mr{जिना} \mr{चढल्यावर} \mr{आजीला} \mr{धाप} \mr{लागली}.
\end{quote}

\begin{quote}
\emph{jinā caḍhlyāvar ājīlā dhāp lāglī.}
\end{quote}

\begin{quote}
Grandmother was out of breath after climbing the stairs.
\end{quote}

Past (to him …): \mr{त्याला} \mr{धाप} \mr{लागली}, \emph{tyālā dhāp lāglī}.

Dative: \mr{मला} \mr{धाप} \mr{लागली}; \mr{धापा} \mr{टाकणे}, to pant.

\begin{cousinweb}
\textbf{\mr{धाप} \mr{लागणे}} (\emph{dhāp lāgṇe}) is a Marathi phrase built on \mr{धाप}, \textquotedblleft{}panting\textquotedblright{}, which CDIAL lists under the expressive \emph{*dhapp} (CDIAL 6729), with \mr{लागणे}.
\end{cousinweb}

\begin{grammarlens}[title={What you've built}]
One more word for describing what you lived and felt.
\end{grammarlens}

\subsection*{Your turn}
\begin{itemize}
\raggedright
  \item \emph{Say it:} \emph{dhāp lāgṇe}
  \item \emph{Say it:} \emph{dhāp lāgṇe}, once more
  \item \emph{Say it:} \emph{jinā caḍhlyāvar ājīlā dhāp lāglī.} — Grandmother was out of breath after climbing the stairs.
  \item \emph{Recall:} say the Marathi for to burp, then the Marathi for to spit, then say \emph{dhāp lāgṇe} again
\end{itemize}

\begin{tcolorbox}[breakable,colback=teal!4,colframe=teal!35!black,title={Before you move on}]
What does \mr{धाप} \mr{लागणे} mean? (\textquotedblleft{}To get breathless\textquotedblright{}.) Say it once more.
\end{tcolorbox}
\section[\texorpdfstring{\mr{ठसका} \mr{लागणे} (ṭhaskā lāgṇe)}{ठसका लागणे (ṭhaskā lāgṇe)}]{\mr{ठसका} \mr{लागणे} (ṭhaskā lāgṇe) — to choke on food}

\label{lesson:MR-C492-thaska-lagne}



\begin{quote}
Before the new one: say the Marathi for to spit, then the Marathi for to get breathless.
\end{quote}

\subsection*{You'll want to know: \mr{ठसका} \mr{लागणे}}
\textbf{\mr{ठसका} \mr{लागणे}} — \emph{ṭhaskā lāgṇe} — \textquotedblleft{}to choke on food\textquotedblright{}.

\begin{quote}
\mr{घाईत} \mr{पाणी} \mr{प्यायल्याने} \mr{त्याला} \mr{ठसका} \mr{लागला}.
\end{quote}

\begin{quote}
\emph{ghāīt pāṇī pyāylyāne tyālā ṭhaskā lāglā.}
\end{quote}

\begin{quote}
He choked from drinking water too fast.
\end{quote}

Past (to him …): \mr{त्याला} \mr{ठसका} \mr{लागला}, \emph{tyālā ṭhaskā lāglā}.

A coughing fit when food or water goes down the wrong way.

\begin{cousinweb}
\textbf{\mr{ठसका} \mr{लागणे}} (\emph{ṭhaskā lāgṇe}) is a Marathi phrase built on \mr{ठसका}, \textquotedblleft{}a choking cough\textquotedblright{}, whose origin is not settled, with \mr{लागणे}.
\end{cousinweb}

\begin{grammarlens}[title={What you've built}]
One more word for describing what you lived and felt.
\end{grammarlens}

\subsection*{Your turn}
\begin{itemize}
\raggedright
  \item \emph{Say it:} \emph{ṭhaskā lāgṇe}
  \item \emph{Say it:} \emph{ṭhaskā lāgṇe}, once more
  \item \emph{Say it:} \emph{ghāīt pāṇī pyāylyāne tyālā ṭhaskā lāglā.} — He choked from drinking water too fast.
  \item \emph{Recall:} say the Marathi for to spit, then the Marathi for to get breathless, then say \emph{ṭhaskā lāgṇe} again
\end{itemize}

\begin{tcolorbox}[breakable,colback=teal!4,colframe=teal!35!black,title={Before you move on}]
What does \mr{ठसका} \mr{लागणे} mean? (\textquotedblleft{}To choke on food\textquotedblright{}.) Say it once more.
\end{tcolorbox}
